% Титульный лист
\thispagestyle{empty}
\begin{center}
    \small \MakeUppercase{\UniversityName}\\
    \small \MakeUppercase{\FacultyName}\\
    \small \MakeUppercase{\DepartmentName}
\end{center}

\vspace{1.5cm}

\begin{center}
    \large \textbf{ОТЧЁТ ПО КУРСОВОЙ РАБОТЕ}\\
    \normalsize за период: \ReportPeriod\\
    \normalsize дата: \ReportDate
\end{center}

\vfill

\begin{center}
    \Large \textbf{\WorkTitle}
\end{center}

\vfill

\noindent
\begin{tabularx}{\textwidth}{l X}
    Выполнил:             & студент группы \StudentGroup{} \StudentName \\
    Научный руководитель: & \SupervisorTitle{} \SupervisorName \\
    Работа выполняется:   & 2026/2027 учебный год (5--6 семестры) \\
    Репозиторий отчётов:  & \texttt{\ReportsRepoUrl} \\
\end{tabularx}

\vfill

\begin{center}
    \CityName, \the\year{} г.
\end{center}

\clearpage

% Аннотация
\section*{Аннотация}

Библиотека \texttt{pysatl-criterion} реализует порядка 235 критериев для 16 семейств
распределений, а \texttt{pysatl-experiment} предоставляет фреймворк Monte-Carlo
экспериментов: сегодня в нём доступны типы \texttt{critical\_value}, \texttt{power},
\texttt{time\_complexity} и \texttt{generation\_only}. При этом асимптотическая эффективность
критериев в проекте не оценивается никак: пользователь может измерить мощность на конечной
выборке и стоимость вычисления, но не может сравнить критерии по скорости убывания ошибки
первого рода при росте объёма данных. Ручной расчёт таких характеристик для десятков
критериев трудоёмок и плохо воспроизводим.

Задача работы~--- разработать численный метод оценки бахадуровской эффективности и
интегрировать его во фреймворк в виде нового типа эксперимента \texttt{bahadur\_efficiency},
переиспользующего выборки, сгенерированные из альтернатив, и уже рассчитанные предельные
распределения статистик при нулевой гипотезе. Эксперимент проходит стандартный конвейер
«генерация~--- исполнение~--- отчёт» и завершается PDF-отчётом. В границы работы не входят
вывод функций больших уклонений для статистик и реализация новых критериев; аналитические
значения эффективности предоставляются коллегой по команде и используются как эталон.

Поскольку большинство критериев библиотеки основано на статистиках с невырожденным предельным распределением,
точный наклон Бахадура для них не определён, поэтому применяется приближённый наклон. Реализация ведётся на
Python $\geq$ 3.10 с использованием NumPy, SciPy, SQLAlchemy; данные хранятся в
SQLite или PostgreSQL, вычисления распараллеливаются. Погрешность оценок количественно оценивается,
корректность подтверждается сопоставлением с аналитическими значениями.

Ожидаемые результаты: новый тип эксперимента, доведённый до состояния, когда он запускается
из CLI и даёт воспроизводимый результат; модуль численной оценки наклона и эффективности с
доверительными интервалами; матрица эффективностей для набора критериев
выбранного семейства по 3--5 альтернативам; пример конфигурации, документация и модульные тесты,
проходящие в CI.

\clearpage
